\documentclass[11pt]{article}
\usepackage[margin=1in]{geometry}
\usepackage{amsmath,amssymb,amsthm}
\usepackage{booktabs}
\usepackage{microtype}
\usepackage[hidelinks]{hyperref}

\newtheorem{lemma}{Lemma}
\newtheorem{theorem}[lemma]{Theorem}
\newtheorem{proposition}[lemma]{Proposition}
\newtheorem{conjecture}[lemma]{Conjecture}
\theoremstyle{definition}
\newtheorem{definition}[lemma]{Definition}

\newcommand{\repourl}{\url{https://github.com/justinlyson/turnpike-solver}}
\newcommand{\Z}{\mathbb{Z}}
\DeclareMathOperator{\supp}{supp}

\title{A Layered, Certified Solver for the Turnpike Problem}
\author{Justin Lyson}
\date{November 2026}

\begin{document}
\maketitle

\begin{abstract}
The turnpike problem asks for $n$ integer points on a line whose $\binom{n}{2}$ pairwise distances form a given
multiset $D$. Backtracking (Skiena, Smith and Lemke, 1990) is fast on typical inputs but exponential on known
families. We describe an open-source solver that runs a sequence of sections, each either polynomial-time or clearly
labelled otherwise: budgeted backtracking, an exact positivity certificate for infeasibility, an anchor-and-square-root
step, factoring of the distance polynomial, factoring after a lift to the lattice of additive relations, a
propagation-and-LP search, and an exhaustive fallback. Every YES answer is verified, and every NO answer comes
from a proof. Sections 1--5 decide part of the input space each. Section~6 decides every input in polynomial time if
one stated conjecture holds. The exhaustive fallback has never been needed on any input we have run, including the
hard families from the literature and inputs built to defeat each section. We give the methods, what is proved,
and measurements.
\end{abstract}

\section{Problem and approach}

Let $X=\{x_1<\dots<x_n\}\subset\Z$ and let $D=\Delta X$ be the multiset of the $\binom n2$ differences $x_j-x_i$,
$i<j$. Given $D$, the \emph{turnpike problem} is to find some $X$ with $\Delta X = D$ (a \emph{realization}) or to
show that none exists. Realizations are unique only up to translation and mirror image, and in general not even then:
distinct \emph{homometric} sets share a distance multiset. Write $W=\max D$; every realization can be shifted
to contain $0$ and $W$.

The classical algorithm~\cite{SSL90} places points from the outside in. The largest distance not yet explained must
join a new point to $0$ or to $W$, so each step has at most two choices. It is fast in practice. Zhang~\cite{Zhang94}
gave inputs on which it takes exponential time. The algebraic view~\cite{RS82,SSL90} encodes $X$ as
$A(z)=\sum_{x\in X}z^x$, so that
\begin{equation}\label{eq:P}
  A(z)\,A(z^{-1}) \;=\; n+\sum_{d} D(d)\,(z^d+z^{-d}) \;=:\; P_D(z),
\end{equation}
and realizations correspond to factors of $P_D$ with $0/1$ coefficients. Factoring is polynomial in the degree $2W$,
so this route is pseudo-polynomial. No worst-case polynomial algorithm is known~\cite{LSS03,Dakic00}.

Our solver combines several methods, each strong where the others are weak, and keeps every answer certified. It runs
the sections below in order and stops at the first that decides (Table~\ref{tab:sections}). The source code, the test
corpus and the benchmarks are at \repourl.

\begin{table}[t]
\centering\small
\begin{tabular}{@{}clll@{}}
\toprule
\# & Section & Decides & Cost \\
\midrule
1 & budgeted backtracking & YES or NO within $n^2$ nodes & polynomial \\
2 & positivity certificate & NO & polynomial \\
3 & anchors and square root & YES & polynomial \\
4 & factoring $P_D$ & YES or NO (complete) & pseudo-polynomial in $W$ \\
5 & lattice lift and factoring & YES or NO (complete) & polynomial in $n,\log W, h$ (Thm.~\ref{thm:t3r}) \\
6 & $T^*_k$ search & YES or NO & polynomial if Conj.~\ref{conj:st} holds \\
7 & exhaustive backtracking & YES or NO & exponential worst case \\
\bottomrule
\end{tabular}
\caption{Sections in run order. A \emph{strict} mode disables every section or option that is not polynomial
(section 4, unprojected lifts, section 7); it may answer UNDECIDED but never wrongly.}
\label{tab:sections}
\end{table}

\begin{proposition}[Soundness]
Every YES answer is verified by recomputing $\Delta X$ exactly and comparing it with $D$. Every NO answer comes from
an exhaustive search, a complete enumeration of factorizations, or an exact certificate. Floating-point values are
used only to choose candidates, never to justify an answer.
\end{proposition}

\section{The sections}

\paragraph{1. Budgeted backtracking.} The rule of~\cite{SSL90} with a budget of $n^2$ nodes (mirror images
are pruned at the root). A run that ends within the budget is exact. Most random inputs finish in about $n$ nodes.

\paragraph{2. Positivity certificate.} By~\eqref{eq:P}, on the unit circle
\[
  f(\theta) \;=\; |A(e^{i\theta})|^2 \;=\; n + 2\sum_d D(d)\cos(d\theta) \;\ge\; 0
\]
for every realizable $D$. Hence one angle with $f(\theta)<0$ proves NO.

\begin{lemma}[Exact angles and aliasing]\label{lem:alias}
For $\theta=2\pi p/q$ with integers $p,q$, $\cos(d\theta)=\cos(2\pi\,(dp \bmod q)/q)$, so the angle reduction is exact
in integer arithmetic for distances of any size. For a grid size $N$, let $a_r=\sum_{d\equiv r \ (N)} D(d)$. Then
$f(2\pi k/N) = n + 2\,\mathrm{Re}\sum_{r} a_r\, e^{2\pi i kr/N}$ for all $k$ simultaneously, which is one FFT of
length $N$ after an $O(|\supp D|)$ fold.
\end{lemma}

The solver evaluates $f$ on all $N$ grid angles: one grid with $N\ge 8W$ for narrow inputs ($8W\le 2^{21}$), and three grids of sizes $2^{18}$,
$3^{11}$, $5^8$ for wide ones, followed by a few random angles. Each candidate is then re-evaluated in ball arithmetic
(Arb, 128 bits); NO is returned only when the whole ball is negative. The cost is polynomial in $n$ and $\log W$. A
positivity failure is only a sufficient condition for NO. In tests it catches nearly every perturbed input at once
(Section~\ref{sec:exp}).

\paragraph{3. Anchors and square root.} Call a point set $\mathcal A\ni 0,W$ of mirror pairs $\{p,W-p\}$ an
\emph{anchor set}. Every other point $q$ satisfies $|q-a|\in D$ for all $a\in\mathcal A$, and $q$, $W-q$ pass this
test together. So the test yields the folded pairs, and what remains unknown is which side of each pair is used. With
$b(z)=\sum_q z^{2q-W}$ over the chosen sides, $b+b^* = S$ (all folded pairs) and $bb^*=P$ (the autocorrelation of the
non-anchor points, known from $D$). Hence $b$ is a root of $t^2-St+P$, that is, $b=(S+\sqrt{S^2-4P})/2$. The square
root of a sparse polynomial is computed exactly, term by term from the top. Anchors are found greedily (a
heuristic). The section is sound in every case, because answers are verified and failure gives UNDECIDED. It decides,
for example, the mirror-symmetric chain of Table~
ef{tab:tricky}, which is slow for backtracking and has a tall lift.

\paragraph{4. Factoring $P_D$.} Factor $z^WP_D(z)$ over $\Z$. Reciprocation $f^*(z)=z^{\deg f}f(1/z)$ permutes the
irreducible factors: a self-reciprocal factor occurs with even multiplicity $2m$, and $A$ contains exactly $f^m$; a
reciprocal pair $f^e(f^*)^e$ contributes $f^j(f^*)^{e-j}$ with $0\le j\le e$~\cite{RS82}. Enumerating these tuples,
up to the mirror symmetry $j\mapsto e-j$, finds all $0/1$ factors, so the section is complete. Candidates are first
screened by their lowest and highest coefficients using truncated products. The cost grows with $W$, so this section
is pseudo-polynomial; it is off in strict mode and limited to $W\le 3000$ by default.

\emph{Towers.} Product sets $A(z)=A_1(z)A_2(z^s)$ are the classical source of homometric sets and of hard
inputs. Their factor $A_2(z^s)$ is found cheaply:

\begin{lemma}[Polyphase gcd]\label{lem:tower}
Write $P(z)=\sum_{r=0}^{s-1}z^rP_r(z^s)$. A polynomial $F(z^s)$ divides $P$ if and only if $F$ divides every
$P_r$. If $\deg A_1<s$, the product $A_1(z)G(z^s)$ has no two terms on the same exponent, so it has $0/1$
coefficients if and only if $A_1$ and $G$ do.
\end{lemma}
\begin{proof}
If $P=F(z^s)G(z)$ with $G=\sum_r z^rG_r(z^s)$, then multiplication by $F(z^s)$ preserves exponents mod $s$, so
$P_r=FG_r$ by uniqueness of the decomposition. Conversely, $F\mid P_r$ for all $r$ gives $F(z^s)\mid P$. In
$A_1(z)G(z^s)$ the term $z^{i+sj}$ determines $(i,j)$ when $0\le i<s$.
\end{proof}

For large degrees the solver takes the largest $s$ with $H=\gcd_rP_r\ne 1$ (each $s$ is screened by one gcd modulo a
prime), writes $P=R(z)H(z^s)$, enumerates the $0/1$ solutions of $R$ and of $H$ separately (recursively), and combines
them. This route only produces verified YES answers. If it finds none, the full factorization runs. On a product of
two hard cores ($n=420$, $\deg P=7904$, $s=93$) it replaces a 17.9\,s factorization, nearly all of it modular, by
gcds and two small factorizations taking about 0.1\,s.

\paragraph{5. Lattice lift.} Let $S=\supp D$ and
\[
  K(D)=\{v\in\Z^S:\ v(a)+v(b)=v(a+b)\ \text{whenever } a,b,a+b\in S\}.
\]
The identity $v(d)=d$ lies in $K(D)$. Take a reduced basis $k_1,\dots,k_r$ of $K(D)$, set
$c(d)=(k_1(d),\dots,k_r(d))\in\Z^r$, and call $h=\max|k_i(d)|$ the height of the lift.

\begin{lemma}[Every realization lifts]\label{lem:lift}
If $Y\ni 0$ realizes $D$, put $L(0)=0$ and $L(y)=c(y)$ for $y\in Y\setminus\{0\}$. Then $L(y')-L(y)=c(y'-y)$ for all
$y<y'$ in $Y$. Hence $A_L(u)=\sum_{y}u^{L(y)}$ is a $0/1$ factor of the Laurent polynomial
$P_L(u)=n+\sum_d D(d)(u^{c(d)}+u^{-c(d)})$ with $P_L=A_LA_L^*$. Since $c$ is injective, every $0/1$ factor decodes to
at most one point set.
\end{lemma}
\begin{proof}
For $0<y<y'$ in $Y$ the values $y$, $y'-y$, $y'$ all lie in $S$ and $y+(y'-y)=y'$, so $c(y')=c(y)+c(y'-y)$. The case
$y=0$ is immediate. Injectivity holds because some integer combination of the coordinates is the identity.
\end{proof}

$K(D)$ is computed exactly. Values with no decomposition $s=a+b$ in $S$ are free parameters, and every other value is an
integer combination of them. The relations are then imposed by iterated integer-kernel refinement, with saturation and
LLL at the end. Inputs whose gaps are ``a large step plus small noise'' have a lift of small height even when $W$ is
huge, and the lift exposes that structure. By default $P_L$ is factored directly in its $r$ variables. In strict mode a
random integer map $\pi:\Z^r\to\Z^3$, injective on $\{\pm c(d)\}$ and never $0$, projects $P_L$ to three variables
first.

\begin{theorem}[Projected lift; internally reviewed]\label{thm:t3r}
With the projection to $3$ variables, section 5 is complete (it finds every realization, so an empty result proves NO)
and runs in time polynomial in $n$, $\log W$ and $h$.
\end{theorem}

Completeness holds because the true $\pi(A_L)$ is among the enumerated factor products. Spurious relations created by
$\pi$ only add candidates, which verification rejects. The height $h$ is not polynomial in general: on the doubling
chain $1,2,4,\dots$ the lift has $h=W=2^L$. Such rigid chains are instead decided by section 1 or section 3.

\paragraph{6. The $T^*_k$ search.} Section 6 walks a tree of candidate classes $c=\min(p,W-p)$ in increasing order, with
two children at each branching class. Every child is tested by exact propagation on a 0/1 face system, by a linear
program whose refutations are confirmed by an exactly checked integer Farkas certificate, and, when two children
survive, by $k$-level probing with the LP inside every leg. The search never backtracks.

\begin{conjecture}[$\mathrm{ST}_k$]\label{conj:st}
For fixed $k$, every child that does not extend to a realization, of a node that does, is refuted by the $T^*_k$ test.
\end{conjecture}

\begin{theorem}\label{thm:st}
If Conjecture~\ref{conj:st} holds for some fixed $k$, the search decides every input after at most $1+2(n-1)$ child
tests, each of polynomial cost, so the turnpike problem is solvable in polynomial time.
\end{theorem}

The steps of this argument other than Conjecture~\ref{conj:st} have written proofs that have been reviewed
independently. For $k=3$ we know of no counterexample. If a wrong child ever survives, the solver detects it (the
subtree dies while a sibling survived), reports it, and continues with section 7.

\paragraph{7. Fallback.} Unbudgeted backtracking. When it runs, the report says so and asks for the input to be
submitted as an issue.

\section{Experiments}\label{sec:exp}

All times are wall-clock times on a desktop PC (Python 3.10, python-flint, HiGHS, numba), in process, and include
the sections that ran before the deciding one.

\paragraph{Corpus.} 97 cases: random sets up to $n=300$ and $W=10^9$; homometric pairs and their products; Zhang's
family; hard narrow cores; rigid chains up to $W=2^{64}$; mirror-symmetric chains up to $W\approx 4\cdot10^{12}$;
mirrored products; and 37 inputs with no solution. All 97 are decided correctly, in 1.8\,s in total. By default the
deciding sections are quick search (70), lift (15), positivity (10) and factoring (2). Strict mode also decides all 97,
using only polynomial sections (quick 70, lift 17, positivity 10). The fallback was never used.

\paragraph{Cross-check.} 16{,}000 random small inputs (solvable, perturbed, narrow and wide products) were each solved
in default and strict mode; 4{,}000 of them were also solved with all solutions requested and with the algebraic
sections alone. Status and complete solution sets match plain exhaustive search in every case.

\paragraph{Hard and large inputs.} Table~\ref{tab:tricky}. Each YES case is followed by no-solution variants (one
distance moved by $1$); all 28 variants are decided by positivity (27) or quick search (1), in at most 1.0\,s each.

\begin{table}[t]
\centering\small
\begin{tabular}{@{}lrrlr@{}}
\toprule
Input & $n$ & $W$ & decided by & time (s) \\
\midrule
\texttt{zhang\_nz15\_W15000} & 77 & 15{,}000 & lift & 0.22 \\
\texttt{zhang\_nz15\_wide} & 77 & $1.0\cdot 10^{9}$ & lift & 0.25 \\
\texttt{zhang\_nz19\_W19000} & 97 & 19{,}000 & lift & 0.21 \\
\texttt{zhang\_nz19\_wide} & 97 & $1.0\cdot 10^{9}$ & lift & 0.32 \\
\texttt{zhang\_nz25\_W25000} & 127 & 25{,}000 & lift & 0.46 \\
\texttt{zhang\_nz25\_wide} & 127 & $1.0\cdot 10^{9}$ & lift & 0.60 \\
\texttt{HM36B4} & 88 & 151 & factor & 0.19 \\
\texttt{HM42B5} & 106 & 183 & factor & 0.29 \\
\texttt{HM48B6} & 124 & 215 & factor & 0.48 \\
\texttt{N21xN20\_narrow} & 420 & 3{,}952 & lift & 1.71 \\
\texttt{N21xN20\_wide} & 420 & $4.2\cdot 10^{7}$ & lift & 1.15 \\
\texttt{HM24B2xN20\_wide} & 1040 & $4.2\cdot 10^{8}$ & lift & 1.55 \\
\texttt{mirprod\_J16} & 73 & $2.1\cdot 10^{7}$ & quick & 0.00 \\
\texttt{symchain\_HM36\_L40} & 172 & $4.4\cdot 10^{12}$ & vieta & 0.54 \\
\texttt{random\_n2000} & 2000 & $1.0\cdot 10^{9}$ & quick & 0.35 \\
\texttt{random\_n5000} & 5000 & $1.0\cdot 10^{12}$ & quick & 3.22 \\
\bottomrule
\end{tabular}
\caption{Hard and large inputs (\texttt{bench/tricky.py}). On the $n=420$ products, plain backtracking had not finished
after $10^7$--$10^8$ nodes. Zhang's family with $nz=13$ takes 58\,s and 708{,}664 nodes for plain backtracking, against
0.15\,s here.}
\label{tab:tricky}
\end{table}

\section{Status}

\begin{center}\small
\begin{tabular}{@{}ll@{}}
\toprule
Statement & Status \\
\midrule
Soundness of every answer & proved (by construction) \\
Positivity certificate, Lemma~\ref{lem:alias}; polyphase gcd, Lemma~\ref{lem:tower} & proved (elementary) \\
Lift completeness, Lemma~\ref{lem:lift} & proved \\
Projected lift polynomial in $n,\log W,h$, Theorem~\ref{thm:t3r} & proved, internally reviewed \\
$\mathrm{ST}_k$ implies polynomial time, Theorem~\ref{thm:st} & proved, internally reviewed \\
Conjecture $\mathrm{ST}_k$ ($k=3$) & open; no counterexample known \\
Fallback never needed & tested on every input above, not proved \\
\bottomrule
\end{tabular}
\end{center}

Speed on benchmarks is evidence of practicality, not a complexity bound. The one missing step to a worst-case
polynomial algorithm is Conjecture~\ref{conj:st}. Any input that reaches section~7 is a counterexample to it, or a
gap in the earlier sections, and we would like to see it.

\begin{thebibliography}{9}
\bibitem{SSL90} S.~S.~Skiena, W.~D.~Smith, P.~Lemke. Reconstructing sets from interpoint distances.
  \emph{Proc.\ 6th ACM Symposium on Computational Geometry}, 332--339, 1990.
\bibitem{LSS03} P.~Lemke, S.~S.~Skiena, W.~D.~Smith. Reconstructing sets from interpoint distances. In
  \emph{Discrete and Computational Geometry: The Goodman--Pollack Festschrift}, 597--631. Springer, 2003.
\bibitem{RS82} J.~Rosenblatt, P.~D.~Seymour. The structure of homometric sets. \emph{SIAM J.\ Algebraic Discrete
  Methods} 3(3):343--350, 1982.
\bibitem{Zhang94} Z.~Zhang. An exponential example for a partial digest mapping algorithm. \emph{J.\ Computational
  Biology} 1(3):235--239, 1994.
\bibitem{Dakic00} T.~Daki\'c. \emph{On the Turnpike Problem}. PhD thesis, Simon Fraser University, 2000.
\end{thebibliography}

\end{document}
